\section{Compilação: JIT especulativo, JIT de carga e AOT}
\label{sec:vms-jit-aot}

\paragraph{HotSpot: especular com perfil.}
A seção~\ref{sec:jit} mostrou o ciclo: interpretar, medir, compilar o que
é quente com C1 e depois C2, especular sobre o que o perfil mostrou e,
se a especulação falha, \emph{desotimizar} e voltar ao interpretador.
O capítulo~\ref{cap:historia-vms} mostra de onde veio esse ciclo.

\paragraph{Graal como JIT.}
A página \citetitle{doc:graalvm-compiler} descreve o Graal como
\enquote{a dynamic compiler, written in Java, that transforms bytecode
into machine code}, integrado à HotSpot. Entre as fases de otimização,
cita \emph{inlining} agressivo, \emph{inlining} polimórfico e uma
\emph{partial escape analysis} que \enquote{can remove the costly
allocations of certain objects}: a escape analysis da seção~\ref{sec:jit},
aplicada também quando o objeto só escapa em alguns caminhos. O
ciclo de especulação e desotimização é o mesmo da HotSpot; muda o
compilador do último nível.

\paragraph{Truffle: o interpretador vira compilador.}
O Graal também é a base do Truffle, que a documentação apresenta como
uma biblioteca para implementar linguagens que \enquote{automatically
deriv[es] high-performance code from interpreters}
(\citetitle{doc:truffle}). O mecanismo é a
\termo{partial evaluation}{Avaliação parcial: especializar um programa
para parte das suas entradas já conhecidas. Aqui, especializar o
interpretador para um programa fixo, o que produz código compilado desse
programa.}: \textcite{paper:truffle-partial-evaluation} definem a
linguagem só escrevendo um interpretador; o código compilado sai
automaticamente por avaliação parcial e, quando uma especulação falha,
\enquote{execution transfers back to the interpreter}. É a desotimização
da HotSpot generalizada para qualquer linguagem.

\paragraph{Native Image: compilar tudo antes.}
\textcite{paper:native-image-initialize-once}, o paper do Native Image,
descrevem o processo: uma aplicação iterativa de análise de
\emph{points-to} e de \emph{snapshot} do heap, seguida de compilação AOT
com um compilador otimizador. Código de inicialização pode rodar no
build, e a execução começa com um heap já populado, o \emph{image heap}
(seção~\ref{sec:jvm-vs-nativo}); é por isso que o trabalho que o Quarkus
adianta para o build (capítulo~\ref{cap:build-time}) rende tanto no nativo. Sem perfil de
execução e sem class loading, não há o que especular em runtime; a
seção~\ref{sec:jvm-vs-nativo} mostrou o efeito no throughput de pico.

\paragraph{BEAM: JIT sem especulação.}
O JIT da BEAM, o BeamAsm, faz outra coisa. A documentação diz que ele
\enquote{provides load-time conversion of Erlang BEAM instructions into
native code on x86-64 and aarch64}: traduz cada módulo quando ele é
carregado, eliminando o custo de despacho do interpretador e
especializando cada instrução pelos tipos dos argumentos. E avisa:
\enquote{BeamAsm does hardly any cross instruction optimizations}
(\citetitle{doc:erlang-beamasm}). Não há perfil, não há inlining
especulativo, não há desotimização. Em troca, o JIT é simples, o código
nativo existe desde a carga e a troca a quente de módulos
(seção~\ref{sec:vms-hot-code}) continua funcionando como antes.

\paragraph{Três respostas para \enquote{quando compilar}.}
\begin{description}
\item[Durante a execução, com perfil] HotSpot e GraalVM JIT: demoram a
  aquecer, chegam ao melhor código de pico, dependem de especulação e
  desotimização.
\item[Na carga, sem perfil] BeamAsm: o código nativo existe desde a
  carga do módulo, com otimização local por instrução.
\item[No build] Native Image: startup rápido e memória menor, ao custo
  do mundo fechado e de compilar sem ver a execução real.
\end{description}
